% What Predicts a Judge's Errors Is Not Inside the Judge
%
% JUDGe 2026 (NeurIPS workshop), full-paper track: 6 pages + references, appendices unlimited.
% Non-archival; later submission to a main venue is permitted.
%
% Written as six pages, not as a compression of the nine-page version. The estimator work that
% occupied that draft is apparatus here: it is what makes the numbers believable, and it lives in
% the appendices with one-line pointers.
%
%   python paper/fill_numbers.py --verify
%   python paper/check_macros.py && python paper/check_refs.py
%   pdflatex main && bibtex main && pdflatex main && pdflatex main
%   python paper/check_pages.py --limit 6
%
% !! Results for judgebench, llmbar and rewardbench2 are mid-rebuild; their macros are PENDING and
% !! check_macros blocks. mmlu and mmlu_pro are complete.
%
% !! ai_use.tex states that full-text confirmation of the bibliography is outstanding. It is.

\documentclass{article}
\usepackage{iclr2027_conference,times}

\usepackage{amsmath,amssymb,amsthm}
\usepackage{booktabs}
\usepackage{enumitem}
\usepackage{microtype}
\usepackage{caption}
\usepackage{graphicx}
\usepackage{hyperref}
\usepackage{url}
\usepackage[capitalise,noabbrev]{cleveref}

\captionsetup{font=small,labelfont=bf,skip=4pt}
\setlength{\abovedisplayskip}{5pt}
\setlength{\belowdisplayskip}{5pt}
\newtheorem{definition}{Definition}

% GENERATED by paper/fill_numbers.py -- do not edit by hand.
% Every quantity printed in main.tex resolves through here.

\newcommand{\AbsenceDeltaDup}{-0.0032}
\newcommand{\AbsenceDeltaNoise}{-0.0035}
\newcommand{\AbsenceDupOnly}{0.876}
\newcommand{\AbsenceNoiseOnly}{0.504}
\newcommand{\AddBlockNarrow}{32}
\newcommand{\AddBlockWide}{4096}
\newcommand{\AlphaFDR}{0.05}
\newcommand{\AnchorDeltaHi}{+0.009}
\newcommand{\AnchorDeltaLo}{-0.010}
\newcommand{\AnchorNSlices}{6}
\newcommand{\AnchorSEHi}{0.009}
\newcommand{\AnchorSELo}{0.005}
\newcommand{\CoefRatioBest}{\textbf{??}}
\newcommand{\ConcatBestBaseline}{0.712}
\newcommand{\ConcatBestMThree}{0.681}
\newcommand{\ConcatBestMTwo}{0.612}
\newcommand{\CoverageAtRiskGain}{\textbf{??}}
\newcommand{\DelongRatioClustered}{2.17}
\newcommand{\DelongSEIndep}{0.0020}
\newcommand{\EqBand}{0.02}
\newcommand{\FamilySize}{30}
\newcommand{\FdrNContrasts}{30}
\newcommand{\FdrNSurvive}{0}
\newcommand{\FloorToClear}{\textbf{??}}
\newcommand{\IdentHi}{0.903}
\newcommand{\IdentLo}{0.610}
\newcommand{\MThreeOnlyBest}{\textbf{??}}
\newcommand{\MTwoOnlyBest}{\textbf{??}}
\newcommand{\MdeFDR}{0.127}
\newcommand{\MdeFDROffset}{\textbf{??}}
\newcommand{\MdeUncorrected}{0.083}
\newcommand{\MedianSE}{0.034}
\newcommand{\MinClassStratum}{5}
\newcommand{\NBoot}{2{,}000}
\newcommand{\NItemsPerSliceMain}{\textbf{??}}
\newcommand{\NJudgesMain}{\textbf{??}}
\newcommand{\NNeededFDR}{2,007}
\newcommand{\NPrimaryTests}{3}
\newcommand{\NumAgreeFpBf}{36.5\%}
\newcommand{\NumBfAcc}{73.5\%}
\newcommand{\NumFpAcc}{50.0\%}
\newcommand{\NumFpBias}{100.0\%}
\newcommand{\OffsetBestBaseline}{\textbf{??}}
\newcommand{\OffsetBestMFour}{\textbf{??}}
\newcommand{\OffsetBestMThree}{\textbf{??}}
\newcommand{\OffsetBestMTwo}{\textbf{??}}
\newcommand{\OffsetLossWide}{0.000}
\newcommand{\PermHi}{0.502}
\newcommand{\PermLo}{0.470}
\newcommand{\PilotSliceN}{200}
\newcommand{\PooledFitGuardWithin}{0.95}
\newcommand{\PosControlMThreeOnly}{\textbf{??}}
\newcommand{\PosControlMTwoOnly}{\textbf{??}}
\newcommand{\SdtAccHigh}{0.943}
\newcommand{\SdtAccLow}{0.505}
\newcommand{\SdtAccMid}{0.802}
\newcommand{\SdtAurocHigh}{0.815}
\newcommand{\SdtAurocLow}{0.466}
\newcommand{\SdtAurocMid}{0.765}
\newcommand{\SdtNullBandHi}{\textbf{??}}
\newcommand{\SdtNullBandLo}{\textbf{??}}
\newcommand{\SdtNullBest}{\textbf{??}}
\newcommand{\SpectralWindow}{4096}
\newcommand{\StratumMinN}{\textbf{??}}
\newcommand{\SwampBaseStrong}{0.882}
\newcommand{\SwampBaseWeak}{0.745}
\newcommand{\SwampConcatWide}{0.567}
\newcommand{\SwampLossNarrow}{0.029}
\newcommand{\SwampLossWeakWide}{0.158}
\newcommand{\SwampLossWide}{0.315}
\newcommand{\TargetEffect}{0.04}
\newcommand{\VerDtypeShift}{0.02\%}
\newcommand{\VerMaxTokens}{3131}
\newcommand{\VerdictDecodBest}{\textbf{??}}
\newcommand{\VerdictDecodPlanted}{1.000}
\newcommand{\VerdictLeakBest}{\textbf{??}}
\newcommand{\VerdictLeakP95}{\textbf{??}}


\title{What Predicts a Judge's Errors\\Is Not Inside the Judge}
\author{Anonymous authors\\Paper under double-blind review}

\begin{document}
\maketitle

\begin{abstract}
When an LLM judge is wrong, what tells you? We compare three sources of signal on the same items,
under a reference that makes them comparable: the judge's own confidence, properties of the item,
and the judge's hidden states. Against a matched first-order observer --- one that decides as well
as the judge but has no insight into its own errors --- the judge's expressed confidence clears the
reference on \NClearConfidence{} of \NSlicesUsable{} evaluated slices. Item difficulty, estimated from other
judges' verdicts, clears it on \NClearBaseline{}. That signal transfers: trained on two judges and
applied to a third never seen in training, it predicts the held-out judge's errors at
\TransferExtLo{}--\TransferExtHi{} conditional AUROC, beating the held-out judge's own confidence on
\TransferExtBeatsMargin{} of \TransferN{} held-out fits. Adding the judge's last-token hidden states
and attention-graph spectral features on top adds nothing measurable, with equivalence tests rather
than failures to reject. The practical reading is that useful signal about judge reliability is
external to the judge and available without access to it, and that white-box access --- the thing a
black-box literature cannot evaluate --- does not appear to add to it at this scale.
\end{abstract}

\section{Introduction}
\label{sec:intro}

LLM judges are used at a scale that precludes checking their verdicts, so a signal that flags
unreliable ones has direct value. Three candidate sources are available, in increasing order of
access required: the judge's own expressed confidence, properties of the item being judged, and the
judge's internal state. The third is the one interpretability research reaches for and the one a
practitioner querying an API cannot obtain.

Comparing them requires a reference, and the conventional one is wrong. Reporting the AUROC of a
correctness probe against $0.5$ implies a comparison to an observer that cannot decide anything. The
appropriate comparison is an observer that decides \emph{as well as this judge} and has no insight
into how it decided; such an observer scores well above $0.5$, and higher as it becomes more
accurate (\cref{sec:floors}). Against that reference the three sources separate cleanly, and they
separate in the opposite direction to the one the interpretability framing suggests.

\paragraph{Contributions.}
\begin{enumerate}[nosep,leftmargin=*]
\item A matched first-order reference for correctness prediction, with a stated validity range
      (\cref{sec:floors}), which makes signals requiring different access comparable on one scale.
\item The comparison itself (\cref{sec:results}): item difficulty clears the reference on
      \NClearBaseline{} of \NSlicesUsable{} slices, the judge's own confidence on \NClearConfidence{}, and
      the difficulty signal transfers to a held-out judge.
\item A white-box negative (\cref{sec:whitebox}): hidden states and attention-graph features add
      nothing measurable on top, reported as equivalence at a pre-specified band.
\end{enumerate}

\section{Related work}
\label{sec:related}

\paragraph{Judge reliability as a document property.} \citet{tan2026conformal} report, on
SummEval, that the width of a conformal prediction set over Likert scores correlates with
judge--human disagreement at $r_s = 0.576$, and that set width agrees across judges, which they read
as capturing document-level difficulty rather than judge-specific noise. That observation is close
in spirit to ours and is independent of it. It is also obtained in a different setting: four
70B-class API judges scoring summarisation quality against human Likert ratings, versus 1.5B--7B
open-weight judges emitting discrete verdicts against benchmark labels here. Whether the two are the
same effect is open, and we do not treat them as convergent evidence. What our setting adds is a
predictive test on a held-out judge, a matched reference against which the size of the effect can be
read, and the white-box comparison their design cannot make.

\paragraph{Metacognition.} Separating first-order sensitivity from insight into one's own
performance is long established; meta-$d'$ \citep{maniscalco2012} expresses the latter in units of
the former, normalised as meta-$d'$/$d'$ \citep{fleming2014}. \Cref{sec:floors} is a Monte-Carlo
form of that reference for the correctness ROC.

\paragraph{Correctness probes.} Reading correctness or truthfulness from hidden states is an
established technique \citep{azaria2023}, as is behavioural self-evaluation \citep{kadavath2022}.
Probing methodology has long warned that a probe's score reflects the probe as well as the
representation \citep{hewitt2019control,belinkov2022}. We add the comparison against a matched
first-order reference, which changes what such a score licenses.

\section{Setup}
\label{sec:setup}

A judge scores an item in one forward pass; the verdict is read from the next-token log-probabilities
of two verdict tokens, and the \emph{margin} is their signed difference. The target is
$y_i = \mathbb{1}[\hat v_i = g_i]$ for judge verdict $\hat v_i$ and label $g_i$. Judges are
\texttt{Qwen2.5-Instruct} at 1.5B, 3B and 7B, one family so that scale is not confounded with
vendor. Banks are MMLU, MMLU-Pro, JudgeBench, LLMBar and RewardBench~2, balanced $50/50$, at
\NItemsPerSliceMain{} items per slice, a slice being one judge $\times$ format.

The headline metric is the \textbf{conditional} AUROC: items are compared only within a
ground-truth stratum, so a feature that merely separates positive from negative items cannot score.
Folds are grouped on normalised question text. Three feature sets are compared:
\textbf{confidence} (margin), \textbf{external} (peer difficulty --- the fraction of \emph{other}
judges answering the item correctly --- plus prompt length, task-start position and subject), and
\textbf{internal} (last-token hidden states; and attention-graph spectral diagnostics, on the
smaller judges). Full definitions, bank construction and the estimator are in
\cref{app:setup,app:estimator}.

\section{The reference, and when it is usable}
\label{sec:floors}

\begin{definition}[First-order reference]
\label{def:sdt}
$F_{\mathrm{SDT}}$ is the distribution of conditional AUROC obtained by applying the estimator to a
signal-detection observer with one latent decision variable, a threshold, and no access to its own
errors, simulated at the judge's own $d'$ and criterion with those re-estimated on each resample.
\end{definition}

A second reference is needed because of a channel specific to this target: within a stratum, $y$ is
determined by the judge's own verdict, so any feature decoding that verdict is a perfect predictor
carrying no self-knowledge. The last-token hidden state decodes it at \VerdictDecodLo{}--%
\VerdictDecodHi{} across every bank and judge here. Pooled estimation blocks it --- such a feature
scores $1$ in one stratum and $0$ in the other --- but only when strata are balanced, and judges are
biased. $F_{\mathrm{V}}$ measures the residual, and a rung must clear
$\phi = \max(Q_{95}(F_{\mathrm{SDT}}), Q_{95}(F_{\mathrm{V}}))$.

\paragraph{Validity range.} $\phi$ is not always usable, and its failures are silent. When a
minority stratum is thin the simulated null is wide (SD up to $0.299$ against a median of $0.052$
here), and its 95th percentile can reach $1.000$, so nothing clears it by construction. We therefore
report $\phi$ only where the null's spread is at most \FloorSdMax{} \emph{and} $\phi \le
\FloorPhiMax{}$; the two conditions catch different failures, a wide null and a saturated one. They
exclude \NSlicesExcluded{} of \NSlicesTotal{} slices in this study, both on the \texttt{single}
format of preference banks where verdict bias leaves one stratum with a handful of items. That
degradation has several distinct presentations and is characterised in \cref{app:thin}.

\section{Results}
\label{sec:results}

\emph{[Three of five banks are mid-rebuild; their macros are pending. MMLU and MMLU-Pro are
complete and the transfer numbers below are from those two.]}

\paragraph{Confidence rarely clears the reference; difficulty usually does.} Across
\NSlicesUsable{} slices with a usable $\phi$, the judge's own confidence exceeds the reference on
\NClearConfidence{}. The baseline combining peer difficulty with item properties exceeds it on
\NClearBaseline{}. Read against $0.5$ the confidence numbers look informative; read against a
matched first-order observer, almost none of them are.

\paragraph{The difficulty signal transfers to an unseen judge.} \Cref{tab:transfer} holds out each
judge in turn, fits on the remaining two, and predicts the held-out judge's errors. The external
predictor uses peer difficulty computed from the training judges only, plus item properties, and
never observes the held-out judge. It attains \TransferExtLo{}--\TransferExtHi{} conditional AUROC
and beats the held-out judge's own confidence on \TransferExtBeatsMargin{} of \TransferN{} fits.
Combining the two is not consistently better than the external predictor alone.

\begin{table}[h]
\centering\small
\caption{Cross-judge transfer. Each row holds out one judge, fits on the other two, and predicts
the held-out judge's errors. \emph{external} never observes the held-out judge.}
\label{tab:transfer}
\begin{tabular}{llccc}
\toprule
bank & held-out judge & external & confidence & $\phi$ \\
\midrule
\multicolumn{5}{l}{\emph{[table body from \texttt{results/transfer\_\_cheap.json} on rebuild]}} \\
\bottomrule
\end{tabular}
\end{table}

\paragraph{A confound we can measure but not resolve.} If the judges resemble one another closely
enough, peer difficulty is a noisy copy of the held-out judge's own correctness, and transfer would
measure inter-judge agreement rather than difficulty generalising. We therefore report agreement
between each held-out judge and its training judges alongside every fit. It spans
\AgreeLo{}--\AgreeHi{} across the panel, a range too narrow to separate the two accounts: with the
confounding variable nearly constant, no amount of data from this panel has leverage on it. Adding
banks does not help, because agreement is a property of the judge pair rather than of the bank, so
more fits estimate the range more precisely without widening it.

What would resolve it is a judge that disagrees more --- a second model family at comparable scale,
which is the same addition the one-family limitation calls for. We state the confound as measured
and open rather than as a hedge, and the agreement range is what a replication should try to widen.

\paragraph{Practical form.} Reported as risk--coverage rather than $\Delta$AUROC, since a $0.02$
AUROC difference has no operational meaning and a coverage change does (\cref{app:coverage}).

\section{Does white-box access add anything?}
\label{sec:whitebox}

This is the question the black-box literature cannot ask, and the reason this is not a re-derivation
of \citet{tan2026conformal}. Adding the judge's last-token hidden states to the external baseline
produces no contrast surviving Benjamini--Hochberg across the run, and two contrasts are
\emph{equivalent} at a pre-specified band of \EqBand{} conditional AUROC --- a measured negative
rather than a failure to reject. Attention-graph spectral features, fitted alone, attain
\MTwoOnlyLo{}--\MTwoOnlyHi{}.

Two caveats bound this. The comparison is an increment over a baseline that already clears $\phi$ on
most slices, so it says internals are redundant with item difficulty, not that they carry nothing.
And the design's detection threshold matters: on the one slice where the increment is neither
significant nor equivalent, resolving an effect of the observed size would need roughly an order of
magnitude more items than any affordable design (\cref{app:power}).

\section{Limitations}
\label{sec:limitations}

One model family and three scales. This is not generic hedging: it is the specific reason the
agreement confound of \cref{sec:results} cannot be resolved here, since three scales of one family
agree over only \AgreeLo{}--\AgreeHi{} and a confound with almost no variation cannot be tested
against. A second family chosen for maximum \emph{disagreement} rather than convenience is the one
addition that would both widen that range and test generality, and it is inference-only on this
arm. Judges are single-token log-probability readouts, which is what
production reward pipelines do but not the only deployment; chain-of-thought judges would need the
generated length as a further nuisance covariate. The external predictor is not free: peer difficulty
requires other judges' verdicts on the same items, so the deployable claim concerns ensembles ---
if several judges are already running, a new one's errors are predictable without running it.
\texttt{single}-format items on preference banks treat a relative preference as absolute truth,
which is where the thin-stratum degradation of \cref{sec:floors} concentrates.

\section{Conclusion}

Measured against an observer that decides as well as the judge but knows nothing about how, the
judge's own confidence is close to uninformative, item difficulty is informative and transfers to
judges never observed, and the judge's internals add nothing detectable on top. The signal is
outside the model.

% ICLR 2027 requires a disclosure of LLM use as a section in the paper text (not counted toward the
% page limit) plus a declaration in the submission form. \input this from main.tex before the
% bibliography.
%
% !! ONE CLAUSE IS NOT YET TRUE. The bibliography paragraph asserts that every reference was
% !! human-verified against the paper it cites. At the time of writing, the references had been
% !! resolved by an LLM reading search results and annotated with what each paper establishes, but
% !! the human confirmation pass had NOT been done. Do that pass, then this section is accurate.
% !! Do not submit with the clause as written until it is. Under the ICLR 2027 AI policy a
% !! misrepresentation of this kind is a Code of Ethics matter, and this paper's own thesis is that
% !! plausible-looking output requires checking.

\section*{Use of large language models}
\label{sec:aiuse}

We disclose LLM involvement in this work in the terms the ICLR 2027 policy requires, and in more
detail than the minimum, because the paper's subject makes the question directly relevant.

\paragraph{As a methodological adversary.} A substantial part of this paper originated in an
extended adversarial exchange with a large language model (Anthropic Claude, Opus~5, accessed
through Claude Code between July and September 2026). The model was given the repository and asked
to attack the estimator rather than to defend it. Several of the findings reported here were first
raised in that exchange and then verified by us in code:

\begin{itemize}[nosep,leftmargin=*]
\item the observation that the ladder's deltas correlated negatively with baseline strength, which
      led to the concat-swamping result of \cref{sec:ladder};
\item the request for a matched signal-detection null, which produced \cref{sec:estimand} and
      retracted a behavioural trend we had been prepared to report as a finding;
\item the identification of the self-verdict leakage channel and the observation that the pooled
      fit, not the conditional metric, is what blocks it (\cref{sec:verdict});
\item the demand for a minimum-detectable-effect calculation before further compute
      (\cref{sec:inference}), which established that our pilot null was uninformative.
\end{itemize}

Every one of these was independently checked: each is reproduced by a planted-data test in the
released test suite, and each is stated in the paper with the measurement rather than the argument.
Where the model's premises were wrong we say so in \cref{sec:corrections}, and two of its
recommendations were tested and rejected on evidence --- notably the proposal to replace the grouped
bootstrap with paired DeLong, which we show would have understated our standard errors by
$\DelongRatioClustered{}\times$.

\paragraph{As an implementer.} The same model wrote or substantially edited code in the released
repository, including the offset ladder estimator, the two null references, the leakage controls,
the power stage and the regression tests. All of it is covered by the test suite and by the build
checks described below; none of it was accepted without a test that fails when the behaviour
regresses.

\paragraph{As a search tool over the literature, with human verification.} Candidate references were
located using LLM-assisted search. Every citation was then opened and checked by a human author
against the claim the sentence makes, and the record of that check --- what each paper establishes,
and which sentence leans on it --- is kept as structured comments in the released
\texttt{refs.bib}. Two mismatches were caught this way and fixed in the \emph{sentence} rather than
the citation: a pair of probing papers that reach opposite prescriptions had been cited as
agreeing, and a reward-model benchmark had been cited behind a later, differently-named dataset it
does not describe. A build check refuses to compile the paper if any cited key lacks a resolvable
DOI or arXiv identifier.

\paragraph{What was not delegated.} The research question, the experimental design, the choice of
banks and panel, the preregistered primary tests, and the decision about what may be claimed are
ours. So is responsibility for everything in this paper, including anything an LLM produced.

\paragraph{Reproducibility of the disclosure itself.} Because a claim about how a paper was written
is as checkable as any other, the repository carries the full provenance: every result file records
the git commit, the package versions --- including the exact VCS commit of the spectral library, not
merely its version string --- the configuration and a digest of the bank that was scored. Every
number printed in this paper is a macro resolved from a results file, and the build fails on an
unresolved one rather than printing a stale value.


\phantomsection\label{sec:refstart}
\bibliographystyle{iclr2027_conference}
\bibliography{refs}

\appendix
\section{Setup details}\label{app:setup}
\section{Estimator, and two ways it fails}\label{app:estimator}
\section{Thin support: one cause, several presentations}\label{app:thin}
\section{Detection threshold and power}\label{app:power}
\section{Risk--coverage}\label{app:coverage}
\section{Errata}\label{app:errata}

\end{document}
